% =====================================================================
%  Lezione 6 — Giochi coalizionali: aspetti computazionali
%  Materiale: Lesson_6_Coalitional_Game_Theory_II.pdf,
%             Lab06-Coalitional_Game_Theory_Computational_Issues.pdf,
%             CoalitionalGamesComputationalIssues_Blank.ipynb
% =====================================================================
\chapter[Giochi coalizionali: aspetti computazionali]{Giochi coalizionali:\\aspetti computazionali}

\parteteoria

La lezione 5 si chiudeva con un problema: la funzione caratteristica ha $2^n$ valori e il
valore di Shapley li usa tutti. Questa lezione mostra tre modi per scalare a molti giocatori.
Il percorso:
\begin{enumerate}[label=\arabic*.]
  \item il problema computazionale;
  \item una classe speciale di giochi: il \kw{gioco dell'aeroporto};
  \item i \kw{giochi su sottografo indotto} (ISG);
  \item una rappresentazione compatta universale: le \kw{MC-Nets};
  \item l'\kw{approssimazione Monte Carlo} del valore di Shapley.
\end{enumerate}

% ---------------------------------------------------------------------
\section{Il problema computazionale}

Un gioco coalizionale è rappresentato da una funzione caratteristica $v: 2^N \to \R$, che ha
un valore per ogni sottoinsieme di giocatori: la sua dimensione è \kw{esponenziale} nel
numero di giocatori. Già con tre giocatori (il gioco del gelato della lezione 5) servono
otto numeri:
\[
\begin{gathered}
  v(\emptyset) = v(\{A\}) = v(\{B\}) = v(\{C\}) = 0,\\
  v(\{A,B\}) = v(\{A,C\}) = 750,\quad v(\{B,C\}) = 0,\quad v(\{A,B,C\}) = 1000.
\end{gathered}
\]
Di conseguenza anche i concetti di soluzione richiedono un numero esponenziale di calcoli.
La formula di Shapley somma su tutte le permutazioni:
\[
  \phi(i, v) = \frac{1}{|N|!} \sum_{\pi \in \Pi_N}
     \Bigl( v\bigl(B(\pi, i) \cup \{i\}\bigr) - v\bigl(B(\pi, i)\bigr) \Bigr).
\]
La versione sui sottoinsiemi scende a $2^{n-1}$ termini per giocatore, ma resta esponenziale.

\paragraph{Tre strade.}
\begin{center}
\renewcommand{\arraystretch}{1.3}
\begin{tabularx}{\linewidth}{@{}l X@{}}
\toprule
\textbf{Strada} & \textbf{Che cosa si ottiene} \\
\midrule
Classi speciali di giochi & formule chiuse per il valore di Shapley, calcolabili in tempo
  polinomiale (gioco dell'aeroporto) \\
Rappresentazione compatta & funzione caratteristica in spazio polinomiale e valore di
  Shapley in tempo polinomiale (ISG, MC-Nets) \\
Approssimazione Monte Carlo & tempo polinomiale per \emph{qualunque} gioco, con garanzie
  sull'errore \\
\bottomrule
\end{tabularx}
\end{center}

% ---------------------------------------------------------------------
\section{Il gioco dell'aeroporto}
\label{sec:aeroporto}

Ci sono $n$ compagnie aeree. Ognuna ha bisogno di una pista di una certa lunghezza per i
propri aerei. Una pista può essere condivisa, quindi conviene costruirne \emph{una sola},
abbastanza lunga per tutti, e dividerne il costo. Come lo si divide in modo equo?

\paragraph{Intuizione.}
La compagnia $A$ ha bisogno di 10~km, $B$ di 15~km, $C$ di 18~km. Si costruisce una pista
di 18~km, ma i suoi tratti servono a gruppi diversi:
\begin{center}
\begin{tikzpicture}[x=0.55cm, y=0.6cm, font=\small]
  \foreach \nome/\km/\y in {A/10/2, B/15/1, C/18/0}{
    \node[anchor=east] at (-0.3,\y) {\nome};
    \draw[-latex, blu, thick] (0,\y) -- (\km,\y) node[right, black] {\km~km};
  }
  \foreach \x in {10,15,18} \draw[dashed, grigio] (\x,-0.6) -- (\x,2.5);
  \node at (5,-1.2)    {$A,B,C$};
  \node at (12.5,-1.2) {$B,C$};
  \node at (16.5,-1.2) {$C$};
  \node at (5,-2.1)    {costo(10 km)$/3$};
  \node at (12.5,-2.1) {costo(5 km)$/2$};
  \node at (16.5,-2.1) {costo(3 km)};
\end{tikzpicture}
\end{center}
I primi 10~km servono a tutti e tre, quindi il loro costo si divide per 3. I 5~km
successivi servono solo a $B$ e $C$, che se li dividono in due. Gli ultimi 3~km servono
solo a $C$, che li paga da solo.

\begin{definizione}[Gioco dell'aeroporto]
I giocatori $N = \{1, \dots, n\}$ hanno costi ordinati $c_1 < c_2 < \dots < c_n$, dove
$c_i$ è il costo della pista che serve al giocatore $i$. La funzione caratteristica è
\[
  v(S) = \max_{i \in S} c_i \qquad \forall S \subseteq N:
\]
una coalizione paga la pista del suo membro più esigente. È un gioco di \kw{costo}: $v(S)$
è una spesa da ripartire, non un guadagno.
\end{definizione}

\begin{teorema}[Littlechild e Owen, 1973]
Il valore di Shapley del gioco dell'aeroporto è
\[
  \phi_i = \sum_{j=1}^{i} \frac{c_j - c_{j-1}}{n - j + 1}
  \qquad \forall i \in N, \qquad c_0 = 0.
\]
\end{teorema}

\paragraph{Lettura della formula.}
Il tratto di pista tra $c_{j-1}$ e $c_j$ costa $c_j - c_{j-1}$ e serve ai giocatori
$j, j+1, \dots, n$, che sono $n - j + 1$: se lo dividono in parti uguali. Il giocatore $i$
paga la sua quota di ogni tratto fino al proprio, cioè per $j = 1, \dots, i$.

\paragraph{Dimostrazione.}
Si definisce $P_j = \{j, \dots, n\}$, l'insieme dei giocatori che hanno bisogno del tratto
$j$, e per ogni $j$ il gioco
\[
  v_j(S) = \begin{cases}
    0 & \text{se } S \cap P_j = \emptyset,\\
    c_j - c_{j-1} & \text{se } S \cap P_j \ne \emptyset.
  \end{cases}
\]
Una coalizione $S$ con membro più esigente $k$ interseca esattamente $P_1, \dots, P_k$,
quindi $\sum_j v_j(S) = \sum_{j=1}^{k} (c_j - c_{j-1}) = c_k = v(S)$: il gioco
dell'aeroporto è la somma dei giochi $v_j$. In $v_j$ i giocatori fuori da $P_j$ sono
\kw{nulli} (non cambiano mai il valore), quelli dentro $P_j$ sono \kw{simmetrici}; per
efficienza si dividono in parti uguali $c_j - c_{j-1}$:
\[
  \phi_i(v_j) = \begin{cases}
    \dfrac{c_j - c_{j-1}}{n - j + 1} & \text{se } i \in P_j \ (\text{cioè } j \le i),\\[2mm]
    0 & \text{altrimenti.}
  \end{cases}
\]
Per \kw{additività}, $\phi_i(v) = \sum_j \phi_i(v_j) = \sum_{j=1}^{i}
\frac{c_j - c_{j-1}}{n-j+1}$. \hfill$\square$

\begin{osservazione}[gli assiomi come strumento di calcolo]
La dimostrazione non tocca mai le $n!$ permutazioni: scompone il gioco in pezzi semplici e
usa i quattro assiomi della lezione 5 (nullo, simmetria, efficienza, additività). È lo
stesso trucco che rende veloci ISG e MC-Nets: \emph{additività} vuol dire che basta
calcolare Shapley su ogni pezzo e sommare.
\end{osservazione}

\begin{esempio}[quattro compagnie, costi $(8, 11, 13, 18)$]
\begin{center}
\renewcommand{\arraystretch}{1.2}
\begin{tabular}{lccccc}
\toprule
 & Tratto 1 & Tratto 2 & Tratto 3 & Tratto 4 & \textbf{Shapley} \\
\midrule
Costo marginale & 8 & 3 & 2 & 5 & \\
\midrule
Quota di P1 & 2 &   &   &   & \textbf{2} \\
Quota di P2 & 2 & 1 &   &   & \textbf{3} \\
Quota di P3 & 2 & 1 & 1 &   & \textbf{4} \\
Quota di P4 & 2 & 1 & 1 & 5 & \textbf{9} \\
\bottomrule
\end{tabular}
\end{center}
Il primo tratto ($8$) si divide fra quattro, il secondo ($11 - 8 = 3$) fra tre, il terzo
($13 - 11 = 2$) fra due, l'ultimo ($18 - 13 = 5$) lo paga solo P4. Totale
$2 + 3 + 4 + 9 = 18 = v(N)$.
\end{esempio}

\paragraph{Complessità.}
La formula come scritta richiede $O(n^2)$ operazioni (una somma di $i$ termini per ogni
giocatore). Ma ogni riga della tabella è la riga precedente più un termine: con una somma
cumulativa si scende a $O(n)$, più $O(n \log n)$ se i costi vanno ordinati. Contro
$O(n \cdot 2^n)$ della formula generale.

\paragraph{Applicazioni.}
Car sharing e ride sharing, capacità di banda in una rete, una condotta d'acqua o un
oleodotto, spazio di archiviazione in cloud con utenti che chiedono quote diverse, livelli
di abbonamento a un software.

\begin{osservazione}[ricetta per riconoscere un gioco dell'aeroporto]
Servono quattro ingredienti:
\begin{itemize}
  \item una sola struttura o risorsa \kw{condivisa};
  \item una struttura \kw{annidata}: chi chiede di più ha bisogno anche di tutto ciò che
        serve a chi chiede di meno;
  \item il costo \kw{cresce} con la richiesta;
  \item l'obiettivo è dividere i costi in modo \kw{equo}.
\end{itemize}
\end{osservazione}

\begin{extra}[il valore di Shapley dell'aeroporto è nel core]
Nel gioco dell'aeroporto una coalizione $S$ con membro più esigente $k$ paga, con Shapley,
$\sum_{i \in S} \phi_i \le$ la somma delle quote di tutti i giocatori fino a $k$, cioè al
più $c_k = v(S)$. Nessun gruppo pagherebbe meno costruendo la pista da solo: la ripartizione
è anche \emph{stabile}. In un gioco di costo la condizione del core ha il verso opposto
rispetto alla lezione 5: $\sum_{i\in S} x_i \le v(S)$. Il gioco è infatti \emph{concavo}
(per i costi), e per questi giochi Shapley sta sempre nel core.
\end{extra}

% ---------------------------------------------------------------------
\section{Giochi su sottografo indotto}
\label{sec:isg}

\begin{definizione}[Induced subgraph game, ISG]
I giocatori sono i nodi di un grafo pesato. Un arco $\{i,j\}$ con peso $w_{ij}$ rappresenta
il valore della coppia; un cappio con peso $w_{ii}$ il valore del giocatore da solo. Il
valore di una coalizione è la somma dei pesi degli archi che hanno \emph{entrambi} gli
estremi nella coalizione, cioè il peso del sottografo indotto da $C$:
\[
  v(C) = \sum_{\{i,j\} \subseteq C} w_{ij}.
\]
\end{definizione}

\begin{teorema}[Shapley negli ISG]
\[
  \phi_i = w_{ii} + \frac12 \sum_{j \in \Gamma(i)} w_{ij},
\]
dove $\Gamma(i)$ è l'insieme dei vicini di $i$.
\end{teorema}

\paragraph{Perché.}
Di nuovo additività: ogni arco $\{i,j\}$ è un gioco a sé, in cui una coalizione vale
$w_{ij}$ se contiene sia $i$ sia $j$, altrimenti 0. Lì $i$ e $j$ sono simmetrici e tutti
gli altri sono nulli, quindi si dividono $w_{ij}$ a metà. Il cappio è un gioco in cui conta
solo $i$, che si prende tutto $w_{ii}$.

\begin{esempio}[ISG con cinque giocatori]
\par\smallskip\noindent
\begin{minipage}{0.4\linewidth}
\centering
\begin{tikzpicture}[x=1cm, y=1cm,
    nodo/.style={circle, draw=blu, fill=blu!15, thick, minimum size=6mm, inner sep=0pt},
    font=\small]
  \node[nodo] (1) at (0,2) {1};
  \node[nodo] (3) at (1.6,1.9) {3};
  \node[nodo] (2) at (0.5,0.8) {2};
  \node[nodo] (4) at (2.2,0.8) {4};
  \node[nodo] (5) at (2.6,-0.5) {5};
  \draw (1) -- node[above]{6} (3);
  \draw (1) -- node[left]{4} (2);
  \draw (2) -- node[above left, pos=0.4]{3} (3);
  \draw (3) -- node[right]{4} (4);
  \draw (2) -- node[below]{3} (4);
  \draw (4) -- node[right]{4} (5);
\end{tikzpicture}
\end{minipage}%
\begin{minipage}{0.6\linewidth}
Valori di alcune coalizioni:
\[
\begin{aligned}
  v(\{1,2,4\}) &= 4 + 3 = 7,\\
  v(\{1,2,5\}) &= 4,\\
  v(\{1,2,3,4\}) &= 4 + 6 + 3 + 4 + 3 = 20.
\end{aligned}
\]
Valori di Shapley (nessun cappio):
\[
\begin{aligned}
  \phi_1 &= \tfrac12(6 + 4) = 5, &
  \phi_2 &= \tfrac12(4 + 3 + 3) = 5,\\
  \phi_3 &= \tfrac12(6 + 3 + 4) = 6{,}5, &
  \phi_4 &= \tfrac12(4 + 3 + 4) = 5{,}5,\\
  \phi_5 &= \tfrac12 \cdot 4 = 2.
\end{aligned}
\]
\end{minipage}
\par\smallskip
Controllo di efficienza: $5 + 5 + 6{,}5 + 5{,}5 + 2 = 24$, la somma di tutti i pesi.
\end{esempio}

\paragraph{Applicazioni.}
Gli ISG modellano reti sociali e relazioni fra giocatori: pesi \kw{positivi} per le
sinergie, pesi \kw{negativi} per i conflitti. Sono anche un modello naturale per la
formazione di coalizioni, che si possono cercare con tecniche di clustering di grafi
(modularità, partizionamento spettrale con il vettore di Fiedler, conduttanza, Infomap,
Walktrap, modelli a blocchi stocastici).

\begin{attenzione}[gli ISG non rappresentano tutti i giochi]
Il valore di una coalizione è determinato dalle sole coppie: non c'è modo di esprimere un
effetto che compare solo quando ci sono \emph{tre} giocatori insieme. Esempio dal notebook:
il gioco dell'aeroporto con costi $(5, 7, 10)$. Singoli e coppie fissano i pesi
($w_{AA} = 5$, $w_{BB} = 7$, $w_{CC} = 10$, $w_{AB} = -5$, $w_{AC} = -5$, $w_{BC} = -7$),
ma allora l'ISG darebbe $v(\{A,B,C\}) = 5$, mentre il valore vero è $10$.
\end{attenzione}

% ---------------------------------------------------------------------
\section{Rappresentazioni compatte}

Una rappresentazione compatta riduce la memoria occupata dalla funzione caratteristica,
di solito usando una struttura di rete. Il prezzo:
\begin{itemize}
  \item il valore di una coalizione non si legge più in $O(1)$ da una tabella, ma si
        \emph{calcola} in tempo polinomiale (sommando pesi);
  \item in cambio, sfruttando l'additività del valore di Shapley, anche Shapley si calcola
        in tempo polinomiale.
\end{itemize}
Gli ISG sono già una rappresentazione compatta: $n$ nodi e al più $n(n+1)/2$ archi (cappi
compresi) invece di $2^n$ valori. Le MC-Nets la generalizzano fino a rappresentare
qualunque gioco.

\subsection{Marginal Contribution Nets}

\begin{definizione}[MC-Net]
Una MC-Net rappresenta la funzione caratteristica come un insieme di \kw{regole}
\[
  \text{pattern} \longrightarrow \text{valore},
\]
dove il pattern è una formula booleana sui giocatori e il valore è il \kw{contributo
marginale} del pattern. Nella forma base il pattern è una congiunzione
$\{a \land b \land \dots \land c\}$, e la regola \kw{si applica} a $S$ se $S$ contiene
tutti i giocatori del pattern. Il valore di una coalizione è la somma dei valori delle
regole che si applicano:
\[
  v(S) = \sum_{\varphi \to x \,:\, \varphi \subseteq S} x.
\]
\end{definizione}

Se i valori possono essere sia positivi sia negativi, già le congiunzioni bastano a
rappresentare \emph{qualunque} gioco.

\begin{esempio}[due regole]
\[
  \{a \land b\} \to 5, \qquad \{b\} \to 2.
\]
\begin{itemize}
  \item $v(\{a\}) = 0$: nessuna regola si applica ($\{a\}$ non contiene $b$);
  \item $v(\{b\}) = 2$: si applica solo la seconda;
  \item $v(\{a, b\}) = 5 + 2 = 7$: si applicano entrambe.
\end{itemize}
\end{esempio}

\paragraph{ISG come MC-Net.}
Ogni arco $\{i,j\}$ di peso $w$ diventa la regola $\{i \land j\} \to w$ e ogni cappio
diventa $\{i\} \to w_{ii}$. Il grafo dell'esempio precedente diventa
\[
  \{1 \land 3\} \to 6,\quad \{1 \land 2\} \to 4,\quad \{2 \land 3\} \to 3,\quad
  \{2 \land 4\} \to 3,\quad \{3 \land 4\} \to 4,\quad \{4 \land 5\} \to 4.
\]
Una MC-Net è in sostanza un \kw{ipergrafo}: ogni regola è un iperarco che collega tutti i
giocatori del pattern, pesato con il suo valore.

\begin{attenzione}[errore nella slide ``Induced subgraph games as MC-Nets'']
La slide scrive $\{4 \land 5\} \to 5$, ma nel grafo l'arco $\{4,5\}$ pesa $4$. Con la regola
della slide verrebbero $\phi_4 = 6$, $\phi_5 = 2{,}5$ e $v(N) = 25$, invece di $5{,}5$, $2$ e
$24$ (verificato al calcolatore).
\end{attenzione}

\begin{teorema}[Shapley nelle MC-Nets con congiunzioni]
\[
  \phi_i = \sum_{\varphi \to x \,:\, i \in \varphi} \frac{x}{|\varphi|}.
\]
\end{teorema}
Ogni regola è un gioco a sé in cui i giocatori del pattern sono simmetrici e gli altri
nulli: il valore $x$ si divide in parti uguali fra i $|\varphi|$ giocatori del pattern. Poi
additività. Nell'esempio a due regole:
\[
  \phi_a = \frac52 = 2{,}5, \qquad \phi_b = \frac52 + \frac21 = 4{,}5,
\]
e $2{,}5 + 4{,}5 = 7 = v(\{a,b\})$.

\begin{extra}[i valori delle regole sono i dividendi di Harsanyi]
Per un gioco qualsiasi i valori delle regole con congiunzioni si ricavano con
l'inclusione-esclusione (trasformata di Möbius):
\[
  d(S) = \sum_{T \subseteq S} (-1)^{|S| - |T|}\, v(T),
\]
e la MC-Net $\{S \to d(S)\}$ rappresenta esattamente $v$. È per questo che le congiunzioni
con valori di segno qualunque sono complete. La formula di Shapley qui sopra dice allora che
ogni giocatore riceve una quota uguale di ogni dividendo delle coalizioni a cui appartiene.

Ma ``completa'' non vuol dire ``compatta''. Per il gioco dell'aeroporto il calcolo dà
$d(S) = (-1)^{|S|+1} \min_{i\in S} c_i$, mai zero: servono tutte le $2^n - 1$ regole
(verificato fino a $n = 10$: 1023 regole). Con costi $(5,7,10)$ si ottiene proprio la MC-Net
del notebook, con valori $5, 7, 10, -5, -5, -7, +5$.
\end{extra}

\subsection{MC-Nets con formule generali}

Le regole possono contenere formule logiche arbitrarie, con letterali \kw{positivi} e
\kw{negativi}. Una regola $\varphi \to x$ si applica a $S$ se $S$ \kw{soddisfa}
$\varphi$, scritto $S \models \varphi$: assegnando ``vero'' ai giocatori in $S$ e ``falso''
a tutti gli altri, la formula risulta vera. Il valore di una coalizione è
\[
  v(S) = \sum_{\varphi \to x \,:\, S \models \varphi} x.
\]

\needspace{14\baselineskip}
\begin{esempio}[MC-Net con un letterale negativo]
\par\smallskip\noindent
\begin{minipage}[t]{0.4\linewidth}
\vspace{0pt}
\[
\begin{aligned}
  x_1 \land x_2 &\to 5\\
  x_2 &\to 2\\
  x_3 &\to 4\\
  x_2 \land \lnot x_3 &\to -2
\end{aligned}
\]
\end{minipage}%
\begin{minipage}[t]{0.6\linewidth}
\vspace{0pt}
\renewcommand{\arraystretch}{1.15}
\begin{tabular}{lll}
\toprule
$S$ & regole applicate & $v(S)$ \\
\midrule
$\{x_1\}$ & nessuna & $0$ \\
$\{x_2\}$ & 2 e 4 & $2 - 2 = 0$ \\
$\{x_3\}$ & 3 & $4$ \\
$\{x_1,x_2\}$ & 1, 2, 4 & $5 + 2 - 2 = 5$ \\
$\{x_1,x_3\}$ & 3 & $4$ \\
$\{x_2,x_3\}$ & 2, 3 & $2 + 4 = 6$ \\
$\{x_1,x_2,x_3\}$ & 1, 2, 3 & $5 + 2 + 4 = 11$ \\
\bottomrule
\end{tabular}
\end{minipage}
\par\smallskip
La regola 4 si applica solo quando c'è $x_2$ \emph{e manca} $x_3$.
\end{esempio}

\begin{teorema}[Shapley nelle MC-Nets generali]
Per una regola $\varphi \to x$ con $p$ letterali positivi e $n$ letterali negativi, ogni
giocatore $i$ che compare in forma positiva e ogni giocatore $j$ che compare in forma
negativa ricevono
\[
  \phi_i = \frac{(p-1)!\,n!}{(p+n)!}\,x = \frac{x}{p\binom{p+n}{n}}, \qquad
  \phi_j = \frac{p!\,(n-1)!}{(p+n)!}\,(-x) = -\frac{x}{n\binom{p+n}{p}}.
\]
Gli altri giocatori sono nulli per quella regola. Per additività il valore di Shapley è la
somma su tutte le regole.
\end{teorema}

\paragraph{Da dove vengono i coefficienti.}
Si guardano solo i $p + n$ giocatori della regola, perché gli altri non contano. Un
giocatore positivo $i$ ``accende'' la regola se arriva \emph{dopo} tutti gli altri positivi
e \emph{prima} di tutti i negativi; la frazione di ordini in cui succede è
$\frac{(p-1)!\,n!}{(p+n)!}$. Un negativo $j$ ``spegne'' la regola (contributo $-x$) se
arriva dopo tutti i positivi e prima degli altri negativi: frazione
$\frac{p!\,(n-1)!}{(p+n)!}$. Senza negativi ($n = 0$) si ritrova $x/p$.

\begin{esempio}[Shapley nella MC-Net con letterale negativo]
\[
\begin{aligned}
  \phi_{x_1} &= \tfrac{1!\,0!}{2!}\cdot 5 = 2{,}5
    &&\text{(solo regola 1)}\\
  \phi_{x_2} &= \tfrac{1!\,0!}{2!}\cdot 5 + 2 + \tfrac{1!\,0!}{2!}\cdot(-2)
    = 2{,}5 + 2 - 1 = 3{,}5 &&\text{(regole 1, 2, 4)}\\
  \phi_{x_3} &= 4 + \tfrac{1!\,0!}{2!}\cdot\bigl(-(-2)\bigr) = 4 + 1 = 5
    &&\text{(regola 3; negativo nella 4)}
\end{aligned}
\]
Somma $2{,}5 + 3{,}5 + 5 = 11 = v(N)$. Il calcolo esponenziale sulla funzione
caratteristica della tabella dà gli stessi valori.
\end{esempio}

\begin{extra}[l'aeroporto con $n$ regole]
Con le formule generali il gioco dell'aeroporto torna compatto. La dimostrazione del
teorema di Littlechild e Owen è già una MC-Net: per ogni tratto $j$ la regola
\[
  (x_j \lor x_{j+1} \lor \dots \lor x_n) \to c_j - c_{j-1}
\]
(``almeno uno di quelli che hanno bisogno del tratto $j$''). Sono $n$ regole invece di
$2^n - 1$; verificato al calcolatore che riproducono $v$. La scelta del linguaggio delle
formule decide quanto è compatta la rappresentazione.
\end{extra}

\paragraph{Riassunto MC-Nets.}
Rappresentazione compatta e \kw{universale}; il valore di una coalizione si calcola in modo
efficiente; il valore di Shapley si calcola in tempo \kw{lineare} nella dimensione
dell'insieme delle regole.

% ---------------------------------------------------------------------
\section{Approssimazione Monte Carlo del valore di Shapley}

\paragraph{L'idea Monte Carlo.}
Se si vuole conoscere una probabilità $P(x)$ ma la si sa solo campionare, la si approssima
con le frequenze: si genera un campione $D$ di dimensione $M$ da $P(X)$ e si stima
\[
  P_D(X = x) = \frac{M_{X=x}}{M},
\]
dove $M_{X=x}$ è il numero di volte in cui è uscito $x$.

\paragraph{Applicata a Shapley.}
Il valore di Shapley è una \emph{media} del contributo marginale su tutte le $n!$
permutazioni, cioè il valore atteso del contributo marginale di $i$ quando la permutazione
è scelta uniformemente a caso. Invece di sommare su tutte, se ne campionano $m$:
\[
  \tilde\phi(i, v) = \frac1m \sum_{\pi \in \mathcal{P}}
     \Bigl( v\bigl(B(\pi, i) \cup \{i\}\bigr) - v\bigl(B(\pi, i)\bigr) \Bigr),
\]
dove $\mathcal{P} \subset \Pi_N$ è un insieme di $m$ permutazioni estratte a caso.

\begin{lstlisting}[style=pseudo]
# Input: v funzione caratteristica; m numero di campioni
def MC_Shapley(v, m):
    $\tilde\phi_i$ = 0 for each i in N
    for k = 1 ... m:
        $\pi_k$ = random permutation of N
        for i = 1 ... n:
            sv = v(B($\pi_k$, i) $\cup$ {i}) - v(B($\pi_k$, i))
            $\tilde\phi_i$ += sv
    for i = 1 ... n:
        $\tilde\phi_i$ = $\tilde\phi_i$ / m
    return $\tilde\phi_1$, ..., $\tilde\phi_n$
\end{lstlisting}

\begin{attenzione}[indice nello pseudocodice delle slide]
Nell'ultimo ciclo le slide scrivono ``\textbf{for} $k = 1 \dots n$'' ma poi dividono
$\tilde\phi_i$: l'indice giusto è $i$, come sopra.
\end{attenzione}

\paragraph{Proprietà.}
\begin{itemize}
  \item \kw{Calcolo efficiente}: $m$ permutazioni, $n$ valutazioni di $v$ ciascuna; se $v$
        si calcola in tempo polinomiale, tutto è polinomiale.
  \item L'\kw{errore} diminuisce come $1/\sqrt{m}$: per dimezzarlo servono quattro volte
        tanti campioni.
  \item Molto efficace in pratica; si applica a \emph{qualunque} gioco e a qualunque
        rappresentazione.
  \item Esistono tecniche di campionamento che riducono la varianza più in fretta:
        \en{antithetic sampling}, estensione multilineare, metodi kernel (riferimento:
        Mitchell et al., 2022, \emph{Sampling Permutations for Shapley Value Estimation},
        JMLR).
\end{itemize}

\begin{extra}[perché $1/\sqrt{m}$, e una garanzia esplicita]
Ogni permutazione dà un contributo marginale $X_k$ indipendente, con media esatta $\phi_i$:
la stima è \emph{non distorta}. La sua varianza è $\sigma_i^2/m$, quindi l'errore tipico
(deviazione standard) è $\sigma_i/\sqrt{m}$. Se i contributi marginali stanno in un
intervallo di ampiezza $r$, la disuguaglianza di Hoeffding dà
\[
  P\bigl(|\tilde\phi_i - \phi_i| \ge \varepsilon\bigr) \le 2\exp\!\left(-\frac{2m\varepsilon^2}{r^2}\right),
\]
cioè con $m \ge \frac{r^2}{2\varepsilon^2}\ln\frac{2}{\delta}$ campioni l'errore è sotto
$\varepsilon$ con probabilità almeno $1 - \delta$. È la ``garanzia sulla soluzione''
promessa all'inizio della lezione; le slide dicono invece che non c'è modo di fissare $m$
esattamente e che in pratica $m \in [1000, 10000]$ funziona bene.
\end{extra}

% ---------------------------------------------------------------------
\section{Riepilogo delle tecniche}

\begin{center}
\renewcommand{\arraystretch}{1.3}
\begin{tabularx}{\linewidth}{@{}>{\raggedright\arraybackslash}p{2.9cm} X X@{}}
\toprule
\textbf{Tecnica} & \textbf{Pro} & \textbf{Contro} \\
\midrule
Classi speciali (aeroporto) & rappresentazione compatta; Shapley veloce &
  espressività limitata \\
Monte Carlo & Shapley veloce; proprietà di convergenza; si applica a qualunque gioco e
  rappresentazione & non esatto; non c'è un modo esatto di scegliere $m$ (in pratica
  $1000$--$10000$) \\
ISG & rappresenta molti giochi; Shapley veloce; poca memoria &
  non completo: non rappresenta qualunque gioco \\
MC-Nets & completo; Shapley veloce; di solito poca memoria &
  memoria $O(2^n)$ nel caso peggiore; sposta il lavoro sul calcolo dei contributi
  marginali, che vanno precalcolati \\
\bottomrule
\end{tabularx}
\end{center}

% =====================================================================
\partelab{Aeroporto, Monte Carlo, ISG e MC-Nets}
% =====================================================================
\fonte{CoalitionalGamesComputationalIssues\_Blank.ipynb, Lab06-Coalitional\_Game\_Theory\_Computational\_Issues.pdf}

\textbf{Obiettivo}: implementare le tre strade della teoria e verificare ogni volta che il
risultato coincide con la formula esponenziale della lezione 5. Il notebook è vuoto (non
ci sono output attesi): tutte le soluzioni qui sotto sono state scritte ed eseguite, e i
numeri riportati sono misurati. Le slide del lab sono un sottoinsieme di quelle della
lezione.

Il notebook riparte dal codice della lezione 5, con una differenza: ora \texttt{powerset}
aggiunge anche \texttt{frozenset()} con valore 0, quindi la formula di Shapley può leggere
\texttt{v[C - \{i\}]} senza \texttt{KeyError} (il problema segnalato nella lezione 5).

\subsection{Il gioco dell'aeroporto}

Il notebook racconta l'origine del problema (Littlechild e Owen, 1973) con un costo di
€1000 al km e piste da 8, 11, 13, 18~km. La funzione caratteristica è già fornita:
\begin{lstlisting}
def airport_game_characteristic_function(cost_list):
    player_list = [str(i) for i in range(len(cost_list))]
    return {subset: (max([cost_list[int(i)] for i in subset])
                     if subset != frozenset() else 0)
            for subset in powerset(player_list)}
\end{lstlisting}
Da scrivere: Shapley con la formula dedicata. La versione lineare accumula la quota del
tratto corrente:
\begin{lstlisting}
def airport_game_shapley_value(cost_list):
    c = sorted(cost_list)
    n = len(c)
    phi, acc, prev = [], 0.0, 0
    for j, cj in enumerate(c):          # restano n - j giocatori
        acc += (cj - prev) / (n - j)
        phi.append(acc)
        prev = cj
    return phi
\end{lstlisting}
Sui costi $(8000, 11000, 13000, 18000)$ restituisce $(2000, 3000, 4000, 9000)$, identico
alla formula esponenziale \texttt{shapley(v)}.

\begin{attenzione}[un ``/2'' di troppo nel notebook]
Nel calcolo passo per passo il notebook scrive per il giocatore 4
``(€18000 $-$ €13000)$/2$ = €5000''. L'ultimo tratto lo paga solo lui: si divide per 1, e
infatti il risultato €5000 è giusto.
\end{attenzione}

\paragraph{Formula dedicata contro formula generale.}
Giochi casuali con costi distinti fra 1 e 99:
\begin{center}
\renewcommand{\arraystretch}{1.2}
\begin{tabular}{rrrrr}
\toprule
$n$ & Coalizioni & Costruire $v$ & Shapley generale & Shapley aeroporto \\
\midrule
12 & 4\,096      & 0,01 s & 0,03 s  & 10 µs \\
16 & 65\,536     & 0,16 s & 0,79 s  & 19 µs \\
20 & 1\,048\,576 & 3,7 s  & 19,3 s  & 16 µs \\
\bottomrule
\end{tabular}
\end{center}
I valori coincidono sempre. Con $n = 20$ il vettore dei costi occupa 216 byte, il
dizionario 42~MB secondo \texttt{sys.getsizeof}.

\begin{attenzione}[\texttt{sys.getsizeof} misura solo il contenitore]
\texttt{getsizeof} non conta gli oggetti contenuti: per il dizionario misura la tabella hash
ma non il milione di \texttt{frozenset} che fanno da chiavi, quindi la memoria vera è molto
più grande di 42~MB. Lo stesso vale più avanti per il grafo: \texttt{getsizeof(G)}
restituisce 48 byte qualunque sia il grafo. Il confronto del notebook va preso come
indicativo.
\end{attenzione}

\subsection{Approssimazione Monte Carlo}

Prima versione, ingenua: si parte dalla funzione caratteristica completa.
\begin{lstlisting}
def permutation_sampling(player_list, m):
    return [random.sample(player_list, len(player_list)) for _ in range(m)]

def predecessors(player, player_list, permutation):
    return frozenset(permutation[:permutation.index(player)])

def shapley_montecarlo(v, m):
    players = list(max(v))
    phi = {p: 0.0 for p in players}
    for perm in permutation_sampling(players, m):
        for p in players:
            B = predecessors(p, players, perm)
            phi[p] += v[B | {p}] - v[B]
    return {p: phi[p] / m for p in players}
\end{lstlisting}
Sull'aeroporto con 20 giocatori e $m = 1000$ impiega 0,09~s (contro 19~s dell'esatto), con
un errore medio percentuale (MAPE) del 9{,}3\%.

\begin{attenzione}[il test del notebook ordina le stime]
La cella di test confronta le stime \emph{ordinate per valore} con i valori esatti. Funziona
solo perché nell'aeroporto, con costi distinti, $\phi_i$ cresce con $i$; ma ordinando si
riassegnano le stime ai giocatori sbagliati quando due sono vicine, e l'errore risulta
\emph{più basso} del vero: sullo stesso campione 6{,}8\% ordinando, 9{,}3\% confrontando
giocatore per giocatore. Il confronto corretto è per nome del giocatore.
\end{attenzione}

\paragraph{Versione migliore.}
Il notebook chiede di non passare la funzione caratteristica completa ($2^n$ valori: già
costruirla costa più che calcolare Shapley esatto con la formula dedicata). Nell'aeroporto i
contributi marginali di una permutazione si calcolano in un'unica passata: chi entra paga
la differenza fra il proprio costo e il massimo raggiunto finora.
\begin{lstlisting}
def airport_mc_fast(cost_list, m):
    n = len(cost_list)
    phi = [0.0] * n
    idx = list(range(n))
    for _ in range(m):
        random.shuffle(idx)
        cur = 0                          # pista costruita finora
        for i in idx:
            if cost_list[i] > cur:
                phi[i] += cost_list[i] - cur
                cur = cost_list[i]
    return [x / m for x in phi]
\end{lstlisting}
Costo $O(m \cdot n)$, senza mai costruire $v$.

\paragraph{Esperimento facoltativo: errore contro $m$.}
MAPE medio su 20 ripetizioni, costi distinti fra 1 e 999:
\begin{center}
\renewcommand{\arraystretch}{1.2}
\begin{tabular}{rrrrr}
\toprule
$m$ & MAPE ($n=20$) & MAPE$\cdot\sqrt{m}$ & MAPE ($n=100$) & Tempo ($n=100$) \\
\midrule
10     & 91{,}1\% & 2{,}88 & 152{,}9\% & 0,3 ms \\
100    & 27{,}0\% & 2{,}70 & 58{,}9\%  & 2,8 ms \\
1\,000  & 8{,}0\%  & 2{,}52 & 19{,}1\%  & 33 ms \\
10\,000 & 2{,}8\%  & 2{,}84 & 6{,}2\%   & 323 ms \\
\bottomrule
\end{tabular}
\end{center}
La pendenza in scala log-log è $-0{,}505$ per $n = 20$ e $-0{,}467$ per $n = 100$:
l'errore scende come $1/\sqrt{m}$, come dice la teoria (MAPE$\cdot\sqrt{m}$ resta
circa costante). Il grafico della slide ``Error reduces as sample increases'' mostra la
stessa retta.

\begin{osservazione}[l'errore relativo pesa sui giocatori piccoli]
Con $n = 20$ e $m = 1000$ il giocatore con la pista più corta ha un errore relativo del 24\%,
quello con la pista più lunga dell'1{,}1\%. Il giocatore più piccolo paga poco (contribuisce
solo quando arriva per primo, cioè in $1/n$ delle permutazioni), quindi il suo contributo è
quasi sempre zero e raramente grande: varianza alta rispetto alla media. Per questo con
$n = 100$ il MAPE è più alto a parità di $m$: ci sono più giocatori con quote piccole.
\end{osservazione}

\subsection{Giochi su sottografo indotto}

Il notebook usa \texttt{networkx}: un grafo casuale di Erdős–Rényi con 10 nodi e 25 archi,
più un cappio per nodo, con pesi interi casuali fra 1 e 100.
\begin{lstlisting}
def ISG_coalition_value(graph, coalition):
    return graph.subgraph(coalition).size(weight='weight')

def ISG_shapley_value(graph):
    phi = {}
    for i in graph.nodes():
        w_ii = graph.edges[i, i]['weight'] if graph.has_edge(i, i) else 0
        phi[i] = w_ii + 0.5 * sum(graph.edges[i, j]['weight']
                                  for j in graph.neighbors(i) if j != i)
    return phi
\end{lstlisting}
Il valore del sottografo indotto è già in \texttt{networkx}: \texttt{subgraph} tiene solo gli
archi interni alla coalizione (cappi compresi) e \texttt{size(weight=...)} ne somma i pesi.
Il confronto con la funzione caratteristica completa (1024 coalizioni) dà valori identici;
la formula ISG impiega 25~µs, quella esponenziale 9~ms. Sull'esempio della teoria si
ritrovano $v(\{1,2,4\}) = 7$, $v(\{1,2,3,4\}) = 20$ e $\phi = (5, 5, 6{,}5, 5{,}5, 2)$.

\begin{attenzione}[con i cappi, \texttt{neighbors(i)} contiene $i$ stesso]
In \texttt{networkx} un nodo con un cappio è vicino di se stesso. Scrivendo la formula
alla lettera, $w_{ii} + \frac12\sum_{j \in \texttt{neighbors}(i)} w_{ij}$, il cappio viene
contato una volta e mezza. Sul grafo del lab la somma dei valori ingenui è 1920,5 invece di
$v(N) = 1633$: l'efficienza fallisce, e la cella di confronto del notebook non stampa il
messaggio ``The Shapley values are the same!''.
Serve il filtro \texttt{if j != i}.
\end{attenzione}

\begin{attenzione}[quanti archi può avere il grafo]
Un commento del notebook dice che un grafo non orientato con $N$ nodi ha al più
$N(N+1)/2$ archi senza cappi, e che con 10 nodi e i cappi si arriva a 65. I conti giusti
sono $N(N-1)/2 = 45$ senza cappi e $N(N+1)/2 = 55$ con i cappi. Il conteggio successivo
``35 archi'' (25 più 10 cappi) è invece corretto.
\end{attenzione}

\subsection{MC-Nets}

Il notebook implementa le MC-Nets nella forma base: una lista di coppie (pattern, valore),
con il pattern come insieme di giocatori in congiunzione. L'esempio è il gioco
dell'aeroporto con costi $(5, 7, 10)$, che come ISG non si può rappresentare:
\begin{lstlisting}
MC_net = [({"A"}, 5), ({"B"}, 7), ({"C"}, 10),
          ({"A", "B"}, -5), ({"A", "C"}, -5), ({"B", "C"}, -7),
          ({"A", "B", "C"}, 5)]

def MC_net_coalition_value(MC_net, coalition):
    return sum(x for P, x in MC_net if set(coalition).issuperset(P))

def MC_net_shapley_value(MC_net):
    phi = {}
    for P, x in MC_net:
        for i in P:
            phi[i] = phi.get(i, 0) + x / len(P)
    return phi
\end{lstlisting}
Risultati: $v(\{A, C\}) = 5 + 10 - 5 = 10$; la funzione caratteristica ricostruita è
quella dell'aeroporto ($v = 10$ per ogni coalizione che contiene $C$, $v(\{A,B\}) = 7$);
il valore di Shapley è $(1{,}67;\ 2{,}67;\ 5{,}67)$ con tutti e tre i metodi (MC-Net,
formula esponenziale, formula dell'aeroporto). L'esempio a due regole della teoria dà
$\phi_a = 2{,}5$, $\phi_b = 4{,}5$.

\begin{attenzione}[\texttt{max(MC\_net)} trova i giocatori solo per caso]
\texttt{from\_MC\_net\_to\_CFG} ricava l'insieme dei giocatori con
\texttt{max(MC\_net)[0]}. Le regole sono tuple (pattern, valore) e le tuple si confrontano
sul primo elemento: per gli insiemi \texttt{>} vuol dire ``sovrainsieme proprio'', come per
\texttt{max(v)} nella lezione 5. Funziona perché la MC-Net del notebook ha una regola con
\emph{tutti} i giocatori. Senza quella regola si perdono giocatori senza alcun errore: con
le regole $\{A,B\} \to 5$, $\{B\} \to 2$, $\{C\} \to 1$ il risultato è $\{A, B\}$ e $C$
sparisce. Il modo robusto è l'unione dei pattern:
\texttt{set().union(*(P for P, \_ in MC\_net))}.
\end{attenzione}

\paragraph{In più: MC-Nets con letterali negativi.}
Una regola diventa una terna (positivi, negativi, valore), e si applica se
\texttt{pos <= S and not (neg \& S)}. Con la formula di Shapley generale della teoria,
sull'esempio $x_1 \land x_2 \to 5$, $x_2 \to 2$, $x_3 \to 4$, $x_2 \land \lnot x_3 \to -2$
si ottiene $(2{,}5;\ 3{,}5;\ 5)$, identico al calcolo esponenziale sulla funzione
caratteristica ricostruita, che coincide con la tabella delle slide.

% ---------------------------------------------------------------------
\begin{riepilogo}
\begin{itemize}
  \item La funzione caratteristica ha $2^n$ valori e Shapley li usa tutti: servono
        scorciatoie.
  \item \textbf{Aeroporto}: $v(S) = \max_{i\in S} c_i$. Shapley:
        $\phi_i = \sum_{j \le i} \frac{c_j - c_{j-1}}{n-j+1}$, ogni tratto di pista diviso
        fra chi lo usa. Si dimostra per additività; si calcola in $O(n)$.
  \item \textbf{ISG}: $v(C)$ è la somma dei pesi degli archi interni a $C$;
        $\phi_i = w_{ii} + \frac12 \sum_{j\in\Gamma(i)} w_{ij}$. Non completo (solo
        interazioni a coppie).
  \item \textbf{MC-Nets}: regole pattern $\to$ valore; $v(S)$ è la somma delle regole
        soddisfatte. Con congiunzioni: $\phi_i = \sum x/|\varphi|$. Con letterali negativi:
        $\frac{(p-1)!\,n!}{(p+n)!}x$ ai positivi, $\frac{p!\,(n-1)!}{(p+n)!}(-x)$ ai
        negativi. Complete, ma $O(2^n)$ regole nel caso peggiore.
  \item \textbf{Monte Carlo}: media dei contributi marginali su $m$ permutazioni casuali;
        stima non distorta, errore $\propto 1/\sqrt{m}$; vale per qualunque gioco.
  \item Il filo comune: \emph{additività} del valore di Shapley (scomporre il gioco in pezzi
        semplici) oppure \emph{campionamento}.
  \item \textbf{Lab}: aeroporto lineare contro esponenziale (16~µs contro 19~s con
        $n=20$), Monte Carlo con pendenza misurata $-0{,}5$, ISG con \texttt{networkx},
        MC-Nets; tutto verificato contro la formula esatta.
\end{itemize}
\end{riepilogo}
